
\setcounter{page}{205}
\begin{proposition}
The following diagram of functors on \(D_{\mathrm{qc}}^+(Y)\) commutes:
\[
\begin{CD}
\Rfst u_* u^! @>{\operatorname{Trf}_u}>> u^* \\
@V{\alpha}VV @| \\
\Rfst g_* g^! @>{\operatorname{Trf}_g}>> u^*
\end{CD}
\]
where the morphism \(\alpha\) is composed of the isomorphism from [II.5.1] and Corollary 6.4.
\end{proposition}

\begin{proof}
Left to the reader. One verifies compatibility by chasing functor definitions. The subtle point: if \(C^{\bullet\bullet}\) is a Cartan--Eilenberg resolution of \(f^*G^\bullet\), then \(v^b(C^{\bullet\bullet})\) need not be a Cartan--Eilenberg resolution of \(v^b f^*G^\bullet\). However, one may find a dominating Cartan--Eilenberg resolution \(D^{\bullet\bullet}\) mapping to \(v^b(C^{\bullet\bullet})\), which suffices for the argument.
\end{proof}

\setcounter{page}{206}
\begin{proposition}
Let \(X,Y\) be locally noetherian preschemes, let \(f\colon X\to\mathbb{P}_Y^n\) be finite, and let \(g\colon\mathbb{P}_Y^n\to Y\) be the projection, with \(gf\) finite. Then there is a commutative diagram on \(D_{\mathrm{qc}}^+(Y)\):
\[
\begin{CD}
\Rfst(gf)_*(gf)^! @>{\operatorname{Trf}_{gf}}>> \id \\
@VVV @| \\
\Rfst g_* g^! @>{\operatorname{Trp}_g}>> \id
\end{CD}
\]
The left vertical arrow is built from [II.5.1] and the residue isomorphism of Proposition 8.2.
\end{proposition}

\begin{proof}
Consider the fibre product \(X\times_Y\mathbb{P}_Y^n\) as in Proposition 8.2. The result follows from Propositions 6.6 and 10.1.
\end{proposition}

\begin{definition}
A morphism \(f\colon X\to Y\) of preschemes is \textit{projectively embeddable} if it factors as \(f = p\circ g\), where \(p\colon\mathbb{P}_Y^n\to Y\) is a projective space projection and \(g\colon X\to\mathbb{P}_Y^n\) is finite.
\end{definition}

\noindent\textbf{Example.}
If \(f\colon X\to Y\) is a projective morphism with \(Y\) quasi-compact and admitting an ample sheaf, then \(f\) is projectively embeddable (EGA II.5.5.4(ii)).

\setcounter{page}{207}
\begin{theorem}
Consider the category \(\operatorname{Lno}\) of locally noetherian preschemes. There exists a unique trace theory for projectively embeddable morphisms \(f\colon X\to Y\), consisting for each such morphism of a functorial morphism
\[
\operatorname{Tr}_f\colon \Rfst f^! \to \id
\]
on \(D_{\mathrm{qc}}^+(Y)\), subject to axioms \(\textsc{TRA 1--TRA 4}\) below.
\end{theorem}

\noindent\textbf{\textsc{TRA 1}}
For a composition \(X\xrightarrow{f}Y\xrightarrow{g}Z\) of projectively embeddable morphisms, the diagram
\[
\begin{CD}
\mathbf{R}(gf)_*(gf)^! @>{\operatorname{Tr}_{gf}}>> \id \\
@VVV @| \\
\Rfst g_*\Rfst f^! @>{\operatorname{Tr}_g\circ\operatorname{Tr}_f}>> \id
\end{CD}
\]
commutes.

\noindent\textbf{\textsc{TRA 2}}
For every finite morphism \(f\), \(\operatorname{Tr}_f\) is compatible with \(d_f\) (the identification \(f^!\cong f^b\)).

\noindent\textbf{\textsc{TRA 3}}
For the projective projection \(f\colon\mathbb{P}_Y^n\to Y\), \(\operatorname{Tr}_f\) is compatible with \(e_f\) (the identification \(f^!\cong f^*\)).

\noindent\textbf{\textsc{TRA 4}}
Let \(f\colon X\to Y\) be projectively embeddable, \(u\colon Y'\to Y\) flat base change, \(X'=X\times_Y Y'\) with projections \(v\colon X'\to X,\ g\colon X'\to Y'\). The diagram
\[
\begin{CD}
u^*\Rfst f^! @>{\operatorname{Tr}_f}>> u^* \\
@V{b_{u,f}}VV @| \\
\Rfst g_* g^! @>{\operatorname{Tr}_g}>> u^*
\end{CD}
\]
commutes.

\setcounter{page}{208}
\begin{proof}
The notation refers to the structure isomorphisms \(c_{f,g},d_f,e_f,b_{u,f}\) from Theorem 8.7.

To define \(\operatorname{Tr}_f\), choose a factorization \(f=p\circ g\) into finite \(g\) and projective projection \(p\). Define \(\operatorname{Tr}_f\) by composing \(c_{g,p},d_g,\operatorname{Trf}_g,e_p,\operatorname{Trp}_p\). Independence of factorization follows by dominating any two factorizations by a third, then applying Proposition 10.4. Axioms \(\textsc{TRA1--TRA4}\) are straightforward diagram chases.
\end{proof}

\setcounter{page}{209}
\begin{remark}
The second main goal of these notes is to extend this trace theory to \textit{arbitrary proper morphisms}. This cannot be done by naive localization, since proper morphisms are not locally projective.

Instead, we abandon the projective case entirely, and construct trace maps for graded residual complexes using only finite morphism traces. After proving the residue theorem, we bootstrap to obtain a general trace map—under the hypotheses that schemes are noetherian of finite Krull dimension admitting dualizing complexes, and complexes have coherent cohomology.
\end{remark}